%% ===================================================================
%%  Towards Longitudinal Prediction of Parkinson's Disease
%%  Motor Subtypes
%%
%%  IEEE journal format (IEEEtran). Figures and table bodies are
%%  generated by `make_assets.py` from committed result files; every
%%  number quoted in the text is re-derived and checked by
%%  `verify_numbers.py`.
%% ===================================================================
\documentclass[journal]{IEEEtran}

\usepackage{amsmath}
\usepackage{amssymb}
\usepackage{graphicx}
\usepackage{booktabs}
\usepackage{multirow}
\usepackage{array}
\usepackage{url}

\graphicspath{{figures/}}

\begin{document}

\title{Towards Longitudinal Prediction of\\
Parkinson's Disease Motor Subtypes}

\author{Parv~Patel,
        Mayank~Tewatia,
        Varsha~Vasudevan,
        and~K.~R.~Sahasranand%
\thanks{P.~Patel, M.~Tewatia and K.~R.~Sahasranand are with the Indian
Institute of Technology Palakkad, Kerala 678623, India.}%
\thanks{V.~Vasudevan is with NeuroComp.}%
\thanks{Data used in the preparation of this article were obtained from the
Parkinson's Progression Markers Initiative (PPMI) database
(\protect\url{www.ppmi-info.org/access-data-specimens/download-data}). For up-to-date
information on the study, visit \protect\url{www.ppmi-info.org}.}}

\markboth{Towards Longitudinal Prediction of Parkinson's Disease Motor Subtypes}%
{Patel \MakeLowercase{\textit{et al.}}: Towards Longitudinal Prediction of PD Motor Subtypes}

\maketitle

\begin{abstract}
The tremor-dominant versus postural-instability/gait-difficulty (PIGD) motor
subtype --- a thresholded ratio of two MDS-UPDRS item means --- is widely used
as a longitudinal endpoint and as a machine-learning target in Parkinson's
disease. We ask whether the predictability of this endpoint is limited by model
capacity or by the reliability of the label itself. Using 1,080 Parkinson's
Progression Markers Initiative participants contributing 3,361 consecutive
labelled visit pairs, we first characterise the endpoint: 35.2\% (95\% CI
31.8--38.8) of observed subtype changes reverse at the very next assessment, a
visit recorded Indeterminate remains so with probability only 0.192
(0.148--0.234), and moving the decision thresholds by less than the spacing
between attainable ratio values relabels 6.7\% of all labelled visits and
21.8\% of participants, every one of them out of the Indeterminate class. We
then benchmark nine specifications spanning five model families --- observed-label
Markov chains, continuous-time hidden Markov models, a hierarchical subscore
mixed model, Gaussian processes on the subscores, and a time-varying spline
model --- under one decision rule, one evaluation set, one patient-clustered
bootstrap and one parameter-free baseline. All nine fall within 1.22 accuracy
points of one another (0.7453--0.7575) and all exceed a persistence rule
(0.7227), but a stratified decomposition shows that 94.3\% of the latent-state
model's margin arises on the 9.9\% of transitions recorded Indeterminate; on
transitions recorded PIGD it gains nothing at all. No fitted model recovers the
Indeterminate class, and the persistence rule attains the highest macro-F1 of
any specification tested. A continuous nine-domain representation of the same
examinations ties with the hidden Markov model on the label yet discriminates
later-recorded Hoehn \& Yahr stage $\geq 3$ and mild cognitive impairment
substantially better than the label does. The binding constraint is the
measurement, not the model.
\end{abstract}

\begin{IEEEkeywords}
Parkinson's disease, disease progression modelling, label noise, benchmarking,
hidden Markov models, Gaussian processes, model calibration, longitudinal data.
\end{IEEEkeywords}

\IEEEpeerreviewmaketitle

% ===================================================================
\section{Introduction}
% ===================================================================

\IEEEPARstart{P}{arkinson's} disease is clinically heterogeneous, and the most
widely used summary of that heterogeneity is the motor subtype. Patients are
classified as tremor dominant (TD), postural-instability/gait-difficulty
(PIGD), or Indeterminate by an arithmetically simple rule: eleven MDS-UPDRS
items are averaged into a tremor score, five into a PIGD score, and the ratio of
the two is compared against a pair of thresholds, with the interval between them
forming the Indeterminate band~\cite{stebbins2013,goetz2008}. Its appeal is that
it is cheap, interpretable, and computable from an examination performed anyway.

That label is used well beyond description. It stratifies cohorts, defines
subgroups for biomarker comparison, and is increasingly treated as a prediction
target for machine-learning models. An international appraisal of Parkinson's
disease subtype classifications examined many such systems and found the
evidence base supporting them thin~\cite{mestre2021,marras2024}; whether the
PIGD construct identifies a distinct phenotype at all has been separately
questioned~\cite{nutt2016,urso2022}.

Several groups have shown the label to be unstable over follow-up. In the de
novo PPMI cohort, substantial fractions of patients changed subtype within
twelve months~\cite{simuni2016}; comparable instability has been reported over
three years in an independent early-disease cohort~\cite{kohat2021}, and
competing subtyping algorithms applied to the same de novo patients agree only
fairly with one another~\cite{ren2020,ren2021}.

What this literature establishes is that the label \emph{changes}. What it does
not establish is \emph{why}, and the distinction matters for anyone who trains a
model on it. Studies of subtype stability measure change as a comparison
between a baseline and a later assessment, or as a shift in prevalence across a
cohort. That design cannot separate a patient whose phenotype genuinely
progressed from one whose single noisy assessment crossed a threshold and
crossed back, because both appear as a changed subtype. Separating them requires
three consecutive assessments and a statistic that asks how often a change is
undone.

The question is sharpened by how predictive studies of this endpoint are
typically reported. Discrimination is reported far more often than
calibration~\cite{vancalster2019}, and we did not find a study of subtype
prediction that reports what a trivial rule --- naming the currently recorded
label --- achieves on the same evaluation set. Without that comparison, an
accuracy figure on a label this persistent is difficult to interpret: a target
whose majority class repeats at better than 80\% rewards a model for doing
nothing.

We therefore asked a question about the measurement rather than about any one
model: \emph{is longitudinal prediction of the motor subtype limited by the
choice of model family, or by the reliability of the label being predicted?}

\subsection{Contributions}

\begin{enumerate}
\item \textbf{A harmonised multi-family benchmark.} Nine specifications spanning
five dissimilar model families are scored on one common set of 3,361 held-out
prediction opportunities, under one decision rule, one metric suite, one
patient-clustered bootstrap and one parameter-free baseline. Prior comparisons
in this literature score different families under different label-derivation
rules, which makes a claim of the form ``model $X$ beats the trivial rule by
$Y$'' ill-defined across the comparison.

\item \textbf{A reliability characterisation of the endpoint itself.} We
quantify immediate reversion of observed changes, the dependence of transition
probabilities on the inter-visit interval, and the sensitivity of the
classification to a threshold perturbation smaller than its own measurement
resolution.

\item \textbf{A stratified margin decomposition.} We localise where a fitted
model beats the trivial rule, by the label recorded at the prediction point.
This turns an aggregate margin into a statement about which cases the model
actually helps on, and in this cohort it collapses the apparent advantage onto
one stratum.

\item \textbf{A dissociation that identifies the ceiling.} A continuous
nine-domain representation of the same examinations ties with the best
subtype-history model on the label while exceeding the label on independent,
later-recorded clinical outcomes --- evidence that the ceiling is a property of
the target, not of the information available.

\item \textbf{A reproducible pipeline.} Every number reported here is re-derived
from committed result files by a checking script shipped with the manuscript
source.
\end{enumerate}

This study proposes no new model. Its unit of contribution is a protocol and a
set of measurements.

% ===================================================================
\section{Related Work}
% ===================================================================

\subsection{Subtype stability}
Simuni \textit{et al.}~\cite{simuni2016} reported the baseline subtype
distribution and one-year instability of the TD/PIGD classification in the de
novo PPMI cohort. Kohat \textit{et al.}~\cite{kohat2021} followed MDS-UPDRS
subtypes over three years in an independent early-disease cohort. Ren
\textit{et al.}~\cite{ren2020,ren2021} compared three published subtyping
algorithms in the same de novo patients and found only fair agreement between
them. These studies measure net drift between assessments. None reports, to our
knowledge, the proportion of observed changes that are reversed at the next
assessment --- the statistic that separates phenotypic progression from
assessment variability.

\subsection{Progression modelling in PPMI}
Severson \textit{et al.}~\cite{severson2021} learned latent disease states from
many clinical measurements in PPMI using an input--output hidden Markov model,
and Fereshtehnejad \textit{et al.}~\cite{fereshtehnejad2017} derived clinical
subtyping criteria with longitudinal validation. Both target \emph{data-driven}
states rather than the published threshold rule. Multistate Markov models that
treat a recorded state as a noisy reading of a latent one, through an explicit
misclassification matrix, are long established in panel-data
analysis~\cite{jackson2011}; we use that construction here as one of five
families rather than as the proposed method.

\subsection{Benchmarking practice}
Van Calster \textit{et al.}~\cite{vancalster2019} document how rarely clinical
prediction models are reported with calibration alongside discrimination. Proper
scoring rules provide the principled alternative to thresholded
accuracy~\cite{gneiting2007}. Our protocol applies both, and adds the element we
found missing throughout this literature: a parameter-free baseline evaluated on
the identical case list, with paired intervals.

% ===================================================================
\section{Materials and Methods}
% ===================================================================

\subsection{Study design}
This was a retrospective secondary analysis of a prospective observational
cohort. The analysis population, the subtype definition, the models and the
validation protocol were all fixed before the harmonised comparison reported
here was run, and every model was scored on one common set of prediction
opportunities.

\subsection{Data source and cohort}
Data came from the Parkinson's Progression Markers Initiative
(PPMI)~\cite{marek2011ppmi}, obtained from the LONI Image and Data Archive under
the PPMI Data Use Agreement. Three sources were used: the PPMI Curated Data Cut
(public release 2026-05-11), which supplies the visit skeleton and covariates;
the item-level MDS-UPDRS Part II (patient questionnaire) and Part III (motor
examination) files, which supply the sixteen ratings from which the subtype is
derived; and the participant status file.

Participants were included if the curated cut assigned them to the sporadic
Parkinson's disease subgroup. Visits were excluded if they were unscheduled or
premature-withdrawal encounters or if they fell beyond fifteen years of
follow-up. Fig.~\ref{fig:flow} gives the full flow with every exclusion counted;
Table~\ref{tab:cohort} gives cohort characteristics.

\begin{figure}[!t]
\centering
\includegraphics[width=\columnwidth]{fig1_flow.pdf}
\caption{Participant flow from the curated data cut to the evaluation set.
Exclusions are counted at each step. The dominant cause of label loss is the
OFF-medication requirement, not incomplete examination: 1,195 of 1,333
unlabelled visits (89.6\%) carried only an ON-medication Part III examination.}
\label{fig:flow}
\end{figure}

\begin{table}[!t]
\centering
\caption{Cohort characteristics. Clinical variables are measured at the
participant's first visit and reported as median (interquartile range). The
subtype distribution at the first visit is over the 1,076 first visits that
carry a label.}
\label{tab:cohort}
\footnotesize
% generated by paper/ieee/make_assets.py -- do not edit
\begin{tabular}{@{}lr@{}}
\toprule
Characteristic & Value \\
\midrule
Participants & 1,080 \\
Patient--visits in cohort & 5,774 \\
Visits with a motor subtype & 4,441 (76.9\%) \\
Male & 709 (65.6\%) \\
Age, y & 64.1 (57.1--70.0) \\
Education, y & 16 (14--18) \\
Disease duration, y & 0.54 (0.30--0.97) \\
MDS-UPDRS Part III (OFF) & 21 (15--28) \\
Hoehn \& Yahr (OFF) & 2.0 (1.0--2.0) \\
MoCA & 27 (26--29) \\
Tremor subscore & 0.55 (0.27--0.82) \\
PIGD subscore & 0.20 (0.00--0.40) \\
Follow-up per participant, y & 3.1 (2.0--6.0) \\
Labelled visits per participant & 3 (2--5) \\
Interval between labelled visits, y & 1.00 (1.00--1.16) \\
\midrule
\multicolumn{2}{@{}l}{\textit{Subtype distribution: first visit / all labelled visits}} \\
\quad Tremor dominant & 772 (71.7\%) / 2,891 (65.1\%) \\
\quad Indeterminate & 109 (10.1\%) / 441 (9.9\%) \\
\quad PIGD & 195 (18.1\%) / 1,109 (25.0\%) \\
\bottomrule
\end{tabular}

\end{table}

\subsection{Subtype definition and label construction}
Motor subtype was computed per visit following Stebbins
\textit{et al.}~\cite{stebbins2013}. The \emph{tremor score} is the mean of
eleven items (2.10; 3.15a--b; 3.16a--b; 3.17a--e; 3.18) and the \emph{PIGD
score} the mean of five (2.12, 2.13, 3.10, 3.11, 3.12). Writing $T$ and $P$ for
the two means and $r = T/P$,
\begin{equation}
\mathrm{subtype} =
\begin{cases}
\mathrm{TD}, & r \geq 1.15 \ \text{ or } (P=0,\ T>0),\\
\mathrm{PIGD}, & r \leq 0.90,\\
\mathrm{Indet.}, & \text{otherwise.}
\end{cases}
\label{eq:rule}
\end{equation}
The ``otherwise'' branch covers $0.90 < r < 1.15$ and the doubly degenerate case
$T = P = 0$. The band is open at both ends, so a ratio falling exactly on a
threshold is assigned to the outer class, and the two degenerate cases are
handled explicitly rather than by division.

Three data-handling rules materially affect the result and are stated in full.
First, PPMI records ``unable to rate'' as the sentinel value 101 \emph{in the
same numeric field as a valid 0--4 rating}; every value outside the valid range
was blanked before any mean was taken. This is not a rounding concern: one such
value among the five PIGD items lifts their mean to roughly 20 and reclassifies
a tremor-dominant visit as PIGD. In the raw Part III file, item 3.12 alone
carries 4,302 such values, of which 3,205 survive into examinations eligible for
this analysis. Second, a subscore was set to missing if \emph{any}
constituent item was missing, matching PPMI's own derivation. Third, the
classification uses the \emph{OFF-medication or untreated} examination, since
levodopa suppresses tremor and scoring ON-state examinations would manufacture
apparent PIGD transitions that are pharmacological rather than pathological.
Where a visit carried more than one eligible Part III examination, the record
with the fewest missing items was retained.

Recomputed labels were validated against PPMI's own two derivations: agreement
was 4,431 of 4,431 (100.00\%) against the curated cut's two-class label, and
1,595 of 1,597 (99.87\%) against a three-class label in an archived 2020 release
covering 390 participants.

\subsection{Missing labels}
Labels were not imputed; analyses used complete cases. Of the 1,333 cohort
visits without a subtype, 1,195 (89.6\%) had only an ON-medication Part III
examination and were therefore excluded by the medication-state rule; 65 had no
Part III record; 57 had an examination that was recorded but incomplete; and 16
had a complete Part III but an incomplete Part II. The dominant mechanism is
therefore treatment-related, and the loss is not random: unlabelled visits were
overwhelmingly from treated participants later in follow-up (median 5.0 years
from baseline against 2.0; 96.8\% on antiparkinsonian medication against
57.9\%). The analysed sample is consequently weighted toward earlier, untreated
follow-up, which attenuates rather than exaggerates observed progression.

\subsection{Prediction task and evaluation protocol}
Visits were ordered within participant by elapsed time from that participant's
first visit. Unlabelled visits were removed before pairing, so a gap year does
not break a sequence; a \emph{transition} is a pair of consecutive \emph{labelled}
visits, carrying the actual elapsed interval between them.

The task is: given a participant's history up to and including visit $k$, and
the elapsed interval $\Delta$ to visit $k+1$, predict the subtype recorded at
visit $k+1$. The model is told \emph{when} the next visit happens, not what is
recorded at it; this is the usual convention for panel data but is an assumption
worth stating, and it means no result here addresses when a patient will next be
seen.

Validation is \textbf{leave-one-participant-out}. For each participant in turn,
every model is refitted from scratch on all other participants and then used to
predict that participant's recorded subtypes one visit at a time. No visit of
the held-out participant enters the fit. This yields 3,361 held-out predictions
from 934 participants for every specification. The continuous-domain
representation of Section~\ref{sec:profile} is instead validated by repeated
patient-level 5-fold cross-validation (4 repeats) because leave-one-out was
measured at $\sim$153\,s per fit; the standardiser, model parameters and all
baselines are refitted inside every fold on training participants only. This
difference between validation schemes is a limitation of the one comparison that
spans them, and is flagged where it arises.

\subsection{Model families}
Five families were compared. All are fitted by maximum likelihood with
deterministic restarts.

\subsubsection{Observed-label Markov chains}
Two chains treating the recorded label as the state. The \emph{discrete-time
chain} estimates a row-stochastic transition matrix from observed pairs,
ignoring the elapsed interval. The \emph{continuous-time chain} estimates a
generator $Q$, with the transition matrix over an interval $t$ given by
$P(t)=\exp(Qt)$; off-diagonal rates are parameterised as $\exp(\theta)$ to keep
them positive.

\subsubsection{Continuous-time hidden Markov model}
A latent phenotype $X(t)$ evolves as a continuous-time Markov chain with
generator $Q$, and what is observed at visit $k$ is not $X(t_k)$ but a noisy
reading $Y_k$ related to it by a misclassification matrix $E$, with
$E_{ab}=\Pr(Y=b \mid X=a)$~\cite{jackson2011}. The likelihood of one
participant's whole sequence marginalises over every latent path by the forward
algorithm. Specifications with three and five latent states were fitted.

A constraint is required. Left unconstrained, the maximum-likelihood solution is
degenerate and identical across random restarts: the latent Indeterminate state
emits uniformly and PIGD becomes absorbing at the parameter boundary. We
therefore require each state to be recorded correctly more often than not,
$E_{aa} \geq 0.5$. The fitted diagonal was 0.941 for TD and 0.842 for PIGD; for
Indeterminate it sat exactly at the floor of 0.500, which is itself a result
about that state rather than about the constraint.

\subsubsection{Hierarchical subscore mixed model}
A joint linear mixed model on the tremor and PIGD subscores with correlated
participant-level random intercepts and slopes, from which the subtype is
derived by applying~\eqref{eq:rule}.

\subsubsection{Gaussian processes on the subscores}
Gaussian-process regression on the two subscore trajectories~\cite{rasmussen2006},
in independent and jointly-covarying forms, again with the subtype derived
by~\eqref{eq:rule}.

\subsubsection{Time-varying spline model}
A linear mixed model whose population rate of change is a spline in time from
baseline, with a constant-rate variant as its nested comparator.

\subsection{Baselines}
The canonical baseline is \textbf{persistence}: name the currently recorded
label. It is the trivial rule a clinician would use, it has no parameters, and
it needs no subscores. Four further reference points are reported and kept
clearly separate from it: subscore last-observation-carried-forward, a
patient-mean subscore, a population-mean subscore, and a marginal null that
always predicts the majority class.

\subsection{Harmonisation}
\label{sec:harmonisation}
The central methodological step is that every model is expressed as a
three-class predictive probability vector over (TD, Indeterminate, PIGD) on the
same 3,361 cases, and the predicted class is the \emph{argmax} of that vector.
The chains and the hidden Markov models are natively on this representation.
The subscore models are converted by Monte Carlo through the project's own
unmodified classification rule --- 10,000 draws per case, with common random
numbers so that the draws cancel in paired contrasts. Halving the number of
draws moved accuracy by at most 0.0015 for any specification, and a vectorised
implementation of~\eqref{eq:rule} was checked against the reference
implementation on 5,000 synthetic cases including degenerate ones
(5,000/5,000 agreement).

Harmonisation is not cosmetic. The same saved predictions from one Gaussian
process score 0.7298 under a point-estimate label derivation and 0.7566 under
the probability-vector derivation --- a difference larger than the entire spread
between the nine fitted specifications.

\subsection{Statistical analysis}
Because one participant contributes several correlated transitions, every
interval reported here is a \textbf{patient-level cluster bootstrap}:
participants are drawn with replacement carrying all of their visits, and every
statistic is recomputed from scratch in each replicate (2,000 draws, seed 42).
Paired contrasts share one resample between the two conditions in every
replicate. For the observed-transition table we compared this against the
binomial Wilson interval and found a design effect above 1 in 8 of 9 cells
(median 1.20, range 0.91--1.46), so the binomial interval is optimistic for most
of that table and the correction is not cosmetic.

Models are compared on accuracy, multiclass Brier score, log loss, macro-F1 and
per-class recall; on the two-class task, additionally on area under the ROC
curve and on calibration. Thresholded accuracy alone is reported nowhere without
a probabilistic companion~\cite{gneiting2007,vancalster2019}. Multiplicity is
largely uncorrected across the specifications compared, and analyses labelled
exploratory were specified after the predictions were generated.

\subsection{Threshold sensitivity}
\label{sec:threshold}
The subtype is a threshold applied to a ratio, so its stability under a small
change to those thresholds is a property of the definition rather than of any
model. We quantified it by rebuilding the entire panel under a perturbed rule
(PIGD if $r<0.91$, TD if $r>1.13$, Indeterminate in between, i.e.\ both
thresholds moved and both band boundaries closed on the opposite side) and
comparing the two label columns. The perturbed rule was specified in the
repository before this analysis and was not selected after inspecting results.

No model is involved and nothing is refitted: the two panels share every visit,
every subscore and every cohort decision and differ only in the label column,
and both invariances were verified rather than assumed. \textbf{No uncertainty
interval is reported for this analysis and none is appropriate} --- it is a
deterministic recomputation on a fixed panel, not an estimate from a sample.
Neither rule is treated as correct; the analysis characterises the label's
fragility, not the correctness of any published labelling.

\subsection{Reproducibility}
The panel is frozen and content-addressed (MD5
\texttt{fb038ddbdc19ab3368e13940f398e7ba}). Each experiment writes to its own
immutable directory with a SHA-256 manifest of its inputs, its code and its
outputs, and carries its own machine-checked validation suite; the runs
underlying this paper passed 15/15, 16/16 and 11/11 checks respectively. Source
code, the asset-generation script and a number-verification script that
re-derives every figure in this text from the committed result files are
available at the repository cited in the Data Availability statement. PPMI data
are not redistributable and must be obtained under the PPMI Data Use Agreement.

% ===================================================================
\section{Results}
% ===================================================================

\subsection{Cohort}
Of 4,788 participants and 19,450 visits in the curated data cut, 1,080
participants contributed 5,774 visits meeting the cohort criteria
(Fig.~\ref{fig:flow}). A motor subtype could be assigned at 4,441 visits
(76.9\%); 934 participants had at least two labelled visits and so contributed
at least one prediction opportunity, giving 3,361 consecutive labelled visit
pairs. At the first visit, 772 participants (71.7\%) were tremor dominant, 109
(10.1\%) Indeterminate and 195 (18.1\%) PIGD --- within one percentage point in
each class of an independent analysis of the same cohort~\cite{simuni2016}. The
median interval between consecutive labelled visits was 1.00 years and the
median follow-up 3.1 years.

\subsection{The recorded label is unstable, and a third of the change reverses}

Across 3,361 transitions, TD was the most persistent state and Indeterminate by
far the least (Fig.~\ref{fig:instability}a, Table~\ref{tab:transitions}). A
visit recorded TD was followed by another TD visit with probability 0.808 (95\%
CI 0.788--0.826), and a PIGD visit by another PIGD visit with probability 0.702
(0.661--0.738). A visit recorded Indeterminate remained Indeterminate with
probability only 0.192 (0.148--0.234), returning to TD in 0.440 (0.381--0.500)
of cases and moving to PIGD in 0.368 (0.318--0.417). PIGD was not absorbing:
29.8\% of PIGD visits were followed by a non-PIGD assessment.

\begin{figure*}[!t]
\centering
\includegraphics[width=\textwidth]{fig2_instability.pdf}
\caption{The endpoint is unstable, and much of the instability is not
progression. (a) One-step transition probabilities over 3,361 consecutive
labelled visit pairs, with patient-level cluster-bootstrap 95\% intervals.
(b) Of every observed change of state, the proportion back at the origin state
one visit later; changes involving the Indeterminate band (gold) reverse far
more often than changes crossing the full width of the scale (blue). (c)
Kaplan--Meier cumulative incidence of conversion to PIGD among the 885
participants not PIGD at their first labelled visit, under two definitions of
the event. Requiring the change to hold across two consecutive labelled visits
halves the estimate.}
\label{fig:instability}
\end{figure*}

\begin{table}[!t]
\centering
\caption{One-step transition probabilities with patient-level cluster-bootstrap
95\% intervals (2,000 draws over 934 participants). Rows sum to one.}
\label{tab:transitions}
\small
% generated by paper/ieee/make_assets.py -- do not edit
\begin{tabular}{@{}llrrl@{}}
\toprule
From & To & $n$ & $\hat{p}$ & 95\% CI \\
\midrule
TD ($n{=}2,261$) & TD & 1,827 & 0.808 & (0.788, 0.826) \\
 & Indet & 181 & 0.080 & (0.068, 0.092) \\
 & PIGD & 253 & 0.112 & (0.098, 0.127) \\
\addlinespace[1.5pt]
Indet ($n{=}334$) & TD & 147 & 0.440 & (0.381, 0.500) \\
 & Indet & 64 & 0.192 & (0.148, 0.234) \\
 & PIGD & 123 & 0.368 & (0.318, 0.417) \\
\addlinespace[1.5pt]
PIGD ($n{=}766$) & TD & 141 & 0.184 & (0.153, 0.218) \\
 & Indet & 87 & 0.114 & (0.091, 0.138) \\
 & PIGD & 538 & 0.702 & (0.661, 0.738) \\
\bottomrule
\end{tabular}

\end{table}

\textbf{A third of observed changes reversed immediately.} Of 673 observed
changes of state, 35.2\% (31.8--38.8) returned to the origin state at the
following visit (Fig.~\ref{fig:instability}b). Reversion was concentrated in
changes involving the Indeterminate band: TD$\to$Indeterminate reversed in 0.560
(0.478--0.642) of cases and PIGD$\to$Indeterminate in 0.500 (0.362--0.648),
while changes crossing the full width of the scale reversed less often ---
TD$\to$PIGD in 0.320 (0.253--0.388) and PIGD$\to$TD in 0.393 (0.305--0.481).

The definition of a conversion therefore dominates its estimated frequency.
Among 885 participants not PIGD at their first labelled visit, 263 converted to
PIGD on a single-assessment definition but only 127 when conversion was required
to hold across two consecutive labelled visits --- a reduction of 51.7\%
(45.7--57.7). Kaplan--Meier cumulative incidence at five years was 35.3\%
(30.8--39.9) under the single-visit definition against 15.3\% (12.0--19.0) under
the sustained one (Fig.~\ref{fig:instability}c).

Transition probabilities also scaled with the interval between visits.
TD$\to$TD fell from 0.820 at intervals of 1.5 years or less to 0.578 beyond four
years, while TD$\to$PIGD rose from 0.100 to 0.289, the two interval-stratified
intervals being non-overlapping. A single pooled matrix therefore averages
one-year and multi-year transitions in whatever proportion attendance happened
to produce, and is not a property of the disease.

\subsection{The Indeterminate band is narrower than the measurement resolution}

Because the tremor score is a mean of eleven integer ratings and the PIGD score
a mean of five, very few ratio values are reachable inside a band 0.25 wide. The
381 Indeterminate visits with a defined ratio take only \textbf{eleven distinct
values}, of which 198 (52.0\%) sit at 0.9091 and 100 (26.2\%) at 1.1364 --- one
hundredth and seven thousandths from a threshold respectively
(Fig.~\ref{fig:margin}b).

Rebuilding the whole panel under the perturbed rule of
Section~\ref{sec:threshold} relabelled \textbf{298 of 4,441 labelled visits
(6.7\%), every one of them from the Indeterminate class}
(Table~\ref{tab:threshold}). No TD or PIGD visit moved. The Indeterminate class
fell from 441 visits to 143, a reduction of 67.6\%, with 198 becoming PIGD and
100 becoming TD --- exactly the two ratio values above, so the entire effect is
two attainable values crossing a line that moved by 0.01 and 0.0064. Measured
over participants rather than visits, the perturbation touched \textbf{235 of
1,080 (21.8\%)}. Agreement with PPMI's own derivations fell from 100.00\% to
97.74\% on the two-class label and from 99.87\% to 92.86\% on the three-class
label.

\begin{table}[!t]
\centering
\caption{Labels under the primary rule against the perturbed rule, over all
4,441 labelled visits. The off-diagonal is confined to one row.}
\label{tab:threshold}
\small
% generated by paper/ieee/make_assets.py -- do not edit
\begin{tabular}{@{}lrrrr@{}}
\toprule
\multirow{2}{*}{Primary rule} & \multicolumn{3}{c}{Perturbed rule} & \multirow{2}{*}{Total} \\
\cmidrule(lr){2-4}
 & TD & Indet & PIGD & \\
\midrule
TD & 2,891 & 0 & 0 & 2,891 \\
Indet & 100 & 143 & 198 & 441 \\
PIGD & 0 & 0 & 1,109 & 1,109 \\
\midrule
Total & 2,991 & 143 & 1,307 & 4,441 \\
\bottomrule
\end{tabular}

\end{table}

A classification that is stable only to changes larger than its own measurement
resolution is not stable in a useful sense. This is a property of the
definition, visible without fitting anything.

\subsection{Nine specifications across five families converge}

Under one decision rule, one evaluation set and one baseline, the nine fitted
specifications fell within \textbf{1.22 accuracy points of one another}, from
0.7453 to 0.7575 (Fig.~\ref{fig:families}, Table~\ref{tab:families}). The
parameter-free persistence rule achieved 0.7227.

\begin{figure*}[!t]
\centering
\includegraphics[width=\textwidth]{fig3_families.pdf}
\caption{Nine specifications drawn from five dissimilar model families (blue)
span 1.22 accuracy points on one common evaluation set, against a
parameter-free persistence rule (red) and four further reference points (grey).
All intervals are paired patient-level cluster bootstraps and are reported in
Table~\ref{tab:families}. The $x$-axis is truncated at 0.60.}
\label{fig:families}
\end{figure*}

\begin{table*}[!t]
\centering
\caption{Harmonised comparison of fourteen specifications on the same 3,361
held-out transitions from 934 participants, under one decision rule (argmax of
the three-class probability vector). $n_{\mathrm{Indet}}$ is the number of times
the specification predicted the Indeterminate class. Paired differences against
the persistence rule use a patient-level cluster bootstrap (2,000 draws).}
\label{tab:families}
\small
% generated by paper/ieee/make_assets.py -- do not edit
\begin{tabular}{@{}llrrrrrr@{}}
\toprule
Specification & Family & Acc. & Brier & Log loss & Macro F1 & $n_{\mathrm{Indet}}$ & $\Delta$ vs.\ persistence (95\% CI) \\
\midrule
\multicolumn{8}{@{}l}{\textit{observed-label chain}} \\
\quad Markov chain, discrete & chain & 0.7474 & 0.390 & 0.707 & 0.4929 & 0 & +0.0247 (+0.0159, +0.0334) \\
\quad CTMC, interval-aware & chain & 0.7453 & 0.393 & 0.710 & 0.4909 & 4 & +0.0226 (+0.0135, +0.0316) \\
\multicolumn{8}{@{}l}{\textit{latent-state HMM}} \\
\quad CT-HMM, 3 latent states & HMM & 0.7539 & 0.369 & 0.673 & 0.5037 & 1 & +0.0312 (+0.0223, +0.0411) \\
\quad CT-HMM, 5 latent states & HMM & 0.7513 & 0.370 & 0.675 & 0.5011 & 0 & +0.0286 (+0.0196, +0.0381) \\
\multicolumn{8}{@{}l}{\textit{subscore LMM}} \\
\quad Joint subscore LMM & LMM & 0.7495 & 0.383 & 0.692 & 0.4977 & 0 & +0.0268 (+0.0136, +0.0398) \\
\multicolumn{8}{@{}l}{\textit{time-varying LMM}} \\
\quad Time-varying LMM, constant & LMM & 0.7495 & 0.383 & 0.692 & 0.4977 & 0 & +0.0268 (+0.0136, +0.0398) \\
\quad Time-varying LMM, spline & LMM & 0.7462 & 0.379 & 0.686 & 0.4937 & 0 & +0.0235 (+0.0102, +0.0364) \\
\multicolumn{8}{@{}l}{\textit{subscore GP}} \\
\quad Subscore GP, independent & GP & 0.7575 & 0.360 & 0.654 & 0.5074 & 0 & +0.0348 (+0.0252, +0.0450) \\
\quad Subscore GP, joint & GP & 0.7563 & 0.359 & 0.652 & 0.5070 & 0 & +0.0336 (+0.0243, +0.0436) \\
\multicolumn{8}{@{}l}{\textit{baseline}} \\
\quad Persistence (subscore LOCF) & baseline & 0.7453 & 0.372 & 0.756 & 0.4970 & 1 & +0.0226 (+0.0146, +0.0311) \\
\quad Patient mean subscore & baseline & 0.7423 & 0.380 & 0.759 & 0.4903 & 3 & +0.0196 (+0.0061, +0.0332) \\
\quad Persistence (recorded label) & baseline & 0.7227 & 0.555 & 7.662 & 0.5559 & 334 & --- \\
\quad Marginal null & baseline & 0.6293 & 0.521 & 0.876 & 0.2575 & 0 & -0.0934 (-0.1202, -0.0691) \\
\quad Population mean subscore & baseline & 0.6254 & 0.517 & 0.872 & 0.3460 & 0 & -0.0973 (-0.1215, -0.0723) \\
\bottomrule
\end{tabular}

\end{table*}

All nine exceeded persistence, by between $+0.0226$ (0.0135--0.0316) and
$+0.0348$ (0.0252--0.0450) accuracy points, every interval excluding zero.
Neither raising the number of latent states nor modelling elapsed time changed
this picture materially: a five-state hidden Markov model scored 0.7513 against
the three-state model's 0.7539, and the interval-aware chain scored 0.7453
against the discrete chain's 0.7474.

\textbf{The five-state figure is not a like-for-like comparison and should not
be read as one.} That model reported convergence in only 404 of 934 folds
(43.3\%), against 934 of 934 for every other specification, and 59.2\% of its
predictions come from folds whose fit did not converge. It is reported because
raising the latent-state count neither helped nor fitted, which is itself the
result. It does not carry the convergence finding: its accuracy is interior to
the range, and dropping it leaves the spread at 1.22 points unchanged.

\textbf{No fitted model recovered the Indeterminate class.} Of fourteen
specifications, five ever predicted Indeterminate at all, and four of those did
so between one and four times in 3,361 opportunities. Only the persistence rule
predicted it with any frequency --- 334 times --- and in consequence the
persistence rule achieved the \textbf{highest macro-F1 of any specification
tested} (0.5559, against 0.5074 for the best fitted model), because it was the
only one to recover any of the 332 Indeterminate outcomes.

\subsection{The advantage over the trivial rule lies in one stratum}

Stratifying the comparison by the label recorded at the prediction point
localised the whole effect (Fig.~\ref{fig:margin}a). For the three-state hidden
Markov model, \textbf{94.3\% of its advantage over persistence came from the 334
transitions (9.9\% of the data) at which the current label was Indeterminate.}
Where the current label was TD it gained six correct predictions out of 2,261;
where the current label was PIGD it gained \textbf{none at all} out of 766, the
two achieving identical accuracy of 0.7023.

\begin{figure*}[!t]
\centering
\includegraphics[width=\textwidth]{fig4_margin.pdf}
\caption{Why the margin over the trivial rule is not what it appears to be.
(a) Extra correct predictions over the persistence rule, decomposed by the
subtype recorded at the prediction point. The discrete Markov chain takes all
83 of its gained cases from the Indeterminate stratum and predicts exactly what
persistence predicts on 100\% of both outer strata. (b) The 381 Indeterminate
visits with a defined ratio take only eleven distinct values; 78\% of them sit
at the two values nearest the thresholds. Moving the thresholds by 0.01 and
0.0064 moves those two values across, and with them 6.7\% of all labelled
visits.}
\label{fig:margin}
\end{figure*}

The two approaches also made the same prediction almost always outside that
stratum: identical calls on 98.4\% of current-TD and 96.9\% of current-PIGD
transitions, against 0.3\% where the current label was Indeterminate. The
discrete Markov chain made the point more starkly still, predicting exactly what
persistence predicted on \textbf{100\%} of both outer strata and taking all 83
of its gained cases from the Indeterminate band. The subscore Gaussian process
was the only family to gain anywhere else, and gained nine and eight cases
respectively.

This localisation is consistent with the rest of the evidence. A chain in the
observed label is rejected by the data --- among transitions whose current label
was TD, the probability of a PIGD outcome rose from 0.086 to 0.393 according to
whether the previous label had been TD or PIGD ($\chi^2 = 156.6$, $\mathrm{df} =
4$, $P = 7.9\times10^{-33}$) --- and adding a latent layer is strongly supported
on information criteria ($\Delta$AIC 224, $\Delta$BIC 186 against the best
alternative on a common joint likelihood). The latent model's fitted mean
sojourn times were 10.1 years in TD, 16.5 in PIGD and 2.1 in Indeterminate.

Taken together, the fitted models did not predict progression better than the
trivial rule. They declined to repeat a class that is rarely repeated, and that
accounts for essentially all of their advantage.

\subsection{Restricting the target to two classes changes the task, not the model}

Excluding the Indeterminate outcomes raises the same model's accuracy from
0.7539 on the three-class task to 0.8371 on a matched two-class evaluation of
3,020 transitions from 898 participants. The increase is largely the removal of
an outcome the model essentially never predicts --- one prediction in 3,361 ---
rather than better prediction.

On this two-class task (Table~\ref{tab:twoclass}) the hidden Markov model and
the persistence rule were \textbf{statistically indistinguishable on accuracy}:
0.8371 (0.8202--0.8540) against 0.8338 (0.8167--0.8509). The continuous
nine-domain representation reached 0.8404 (0.8248--0.8566), again
indistinguishable from the hidden Markov model (paired difference $+0.0033$,
$-0.0050$ to $+0.0112$).

\begin{table*}[!t]
\centering
\caption{Two-class (TD vs.\ PIGD) prediction on 3,020 matched transitions from
898 participants. Intervals are patient-level cluster bootstraps.
$^{\dagger}$The persistence rule emits only zero and one, so its log loss is
clipped and it cannot be calibrated at all.}
\label{tab:twoclass}
\small
% generated by paper/ieee/make_assets.py -- do not edit
\begin{tabular}{@{}lrlrrl@{}}
\toprule
Model & Acc. & 95\% CI & Brier & Log loss & AUC (95\% CI) \\
\midrule
Nine-domain profile & 0.8404 & (0.8248, 0.8566) & 0.1177 & 0.3766 & 0.8902 (0.8748, 0.9053) \\
CT-HMM (3 latent states) & 0.8371 & (0.8202, 0.8540) & 0.1215 & 0.3963 & 0.8583 (0.8411, 0.8755) \\
Markov chain, discrete & 0.8298 & (0.8130, 0.8471) & 0.1327 & 0.4311 & --- \\
Persistence rule & 0.8338 & (0.8167, 0.8509) & 0.1598 & 4.1959$^{\dagger}$ & 0.7927 (0.7725, 0.8127) \\
Last observation carried forward & 0.8252 & (0.8083, 0.8425) & 0.1235 & 0.3941 & --- \\
Population mean trajectory & 0.6914 & (0.6677, 0.7153) & 0.2083 & 0.6067 & --- \\
\bottomrule
\end{tabular}

\end{table*}

The models were nonetheless separable from the trivial rule, and from each
other, on measures that use the probabilities rather than the thresholded call.
The hidden Markov model's two-class Brier score was 0.1215 against persistence's
0.1598, and its predicted probabilities tracked observed frequency to within
0.022 on average across six bins spanning the probability range; persistence
emits only zero and one and cannot be calibrated at all. The nine-domain
representation discriminated better than the hidden Markov model by area under
the curve, 0.8902 (0.8748--0.9053) against 0.8583 (0.8411--0.8755), a paired
difference of $+0.0319$ (0.0229--0.0419). We note that this one comparison
spans two validation schemes, which favours the hidden Markov model.

\subsection{Neither more history nor matched retraining improves prediction}

Restricting the hidden Markov model's prediction-time history to the two most
recent visits left performance indistinguishable from the full record on every
metric: accuracy $-0.0041$ ($-0.0189$ to $+0.0100$), balanced accuracy
$-0.0063$, AUC $-0.0045$, Brier $+0.0051$ and log loss $+0.0112$, all five
intervals containing zero. A single visit was measurably worse --- accuracy
$-0.0218$ ($-0.0423$ to $-0.0014$) and AUC $-0.0617$ ($-0.0858$ to $-0.0385$) ---
so the second visit carries real information while the fourth and beyond do not.

Refitting the model on two-visit sliding windows, so that training matched the
restricted input, did not help and degraded the probabilities: balanced accuracy
$-0.0312$ ($-0.0441$ to $-0.0183$) against the full-sequence model. A $2\times2$
design separated the two factors and attributed the loss to the change in
training ($-0.0248$, $-0.0354$ to $-0.0147$, with prediction history held fixed)
rather than to the history restriction ($-0.0064$, interval containing zero).

Prediction was also no more accurate later in follow-up than at the first
opportunity. Predicting the second visit from the first gave 0.8347 for the
hidden Markov model and 0.8455 for the nine-domain representation; no change
from stage one excluded zero for either model at any stage with a reliable
sample size. Balanced accuracy did rise with follow-up, but so did that of the
persistence rule, and the PIGD share of outcomes rose from 23\% to 48\% over the
same stages, so the trend reflects cohort composition rather than anything the
models learn.

\subsection{A continuous representation carries information the label does not}
\label{sec:profile}

The nine-domain representation gives every participant, in each of nine clinical
domains (bradykinesia/rigidity, tremor, PIGD, MDS-UPDRS Parts I and II, MoCA,
GDS, SCOPA-AUT and ESS), two numbers filtered from their own history by a
continuous-time local linear trend model: a \emph{level} and a \emph{rate of
change}. Levodopa-equivalent daily dose enters as an input to the observation
equation rather than as a domain, so that ``has started treatment'' is not one
of the coordinates the representation is built from. Inference is a Kalman
filter, so the state after visit $k$ is a function of visits $0\ldots k$ and
nothing else; ``no future information'' is a structural property of the
recursion rather than a convention to be policed.

The representation was genuinely multidimensional rather than a severity score
in disguise: seven principal components were required for 90\% of its
across-participant variance, with the first explaining 42.2\% of the levels and
37.2\% of the rates (Fig.~\ref{fig:outcomes}a).

\begin{figure*}[!t]
\centering
\includegraphics[width=\textwidth]{fig5_outcomes.pdf}
\caption{The continuous representation carries prognostic information the
categorical label does not. (a) Cumulative variance explained across
participants; seven of nine components are needed for 90\%. (b) Discrimination
of outcomes recorded later and never used in fitting. $^{*}$ marks a paired
difference against the subtype label whose 95\% interval excludes zero. Every
coordinate was selected on the same data against which it was compared, so each
difference is an upper bound.}
\label{fig:outcomes}
\end{figure*}

Against outcomes recorded later and never used in fitting, single coordinates of
that representation discriminated better than the subtype label
(Fig.~\ref{fig:outcomes}b, Table~\ref{tab:outcomes}). For Hoehn \& Yahr stage 3
or worse, the best level coordinate achieved an AUC of 0.685 against the label's
0.579, a difference of $+0.105$ (0.054--0.156); for mild cognitive impairment
derived from verbal-learning and visuospatial testing rather than from the MoCA,
0.710 against 0.549, a difference of $+0.161$ (0.109--0.213). For motor
complications, an outcome the label itself discriminates at chance (0.504), the
difference was $+0.027$ ($-0.039$ to $+0.094$).

\begin{table*}[!t]
\centering
\caption{Discrimination of later-recorded outcomes by a single coordinate of the
nine-domain representation against the subtype label. Coordinates were selected
on the same data against which they were compared, so each $\Delta$ is an upper
bound. $^{\ddagger}$The investigator cognitive diagnosis is not fully
independent of the feature set and is reported for completeness only.}
\label{tab:outcomes}
\small
% generated by paper/ieee/make_assets.py -- do not edit
\begin{tabular}{@{}llrrrrl@{}}
\toprule
Outcome & Best level coordinate & $n$ & Events & Label AUC & Coord.\ AUC & $\Delta$ (95\% CI) \\
\midrule
Hoehn \& Yahr $\geq 3$ & \texttt{updrs2} & 761 & 127 & 0.579 & 0.685 & +0.105 (+0.054, +0.156) \\
MCI (HVLT/Benton) & \texttt{moca} & 818 & 298 & 0.549 & 0.710 & +0.161 (+0.109, +0.213) \\
Cognitive diagnosis$^{\ddagger}$ & \texttt{moca} & 816 & 327 & 0.558 & 0.678 & +0.119 (+0.069, +0.169) \\
Motor complications & \texttt{brady\_rigid} & 736 & 470 & 0.504 & 0.531 & +0.027 (-0.039, +0.094) \\
\midrule
\multicolumn{7}{@{}l}{\textit{Best rate-of-change coordinate, same outcomes}} \\
\quad Hoehn \& Yahr $\geq 3$ & \texttt{updrs2} & 761 & 127 & 0.579 & 0.633 & +0.053 (-0.013, +0.119) \\
\quad MCI (HVLT/Benton) & \texttt{moca} & 818 & 298 & 0.549 & 0.594 & +0.045 (-0.005, +0.096) \\
\quad Cognitive diagnosis$^{\ddagger}$ & \texttt{updrs2} & 816 & 327 & 0.558 & 0.594 & +0.036 (-0.013, +0.085) \\
\quad Motor complications & \texttt{updrs1} & 736 & 470 & 0.504 & 0.537 & +0.033 (-0.026, +0.090) \\
\bottomrule
\end{tabular}

\end{table*}

\textbf{The rates were not established as informative on any outcome.} The best
rate coordinates gave differences of $+0.053$, $+0.045$, $+0.033$ and $+0.036$,
every interval including zero. The same pattern appeared on the representation's
own primary target: predicting the next visit's nine measurements improved on
carrying the last observation forward by $+0.0733$ log-density per domain
observation (0.0657--0.0813), of which the rate-of-change component supplied
$+0.0040$ (0.0019--0.0062) and the medication input $+0.0058$ (0.0041--0.0075)
--- together under a sixth. Most of the gain was the level being filtered across
a participant's history rather than read off the most recent visit.

\subsection{Exploratory: performance varies by participant, but so does the
trivial rule's}

Weighting each participant equally rather than each visit changed the headline
little --- mean per-participant accuracy 0.834 (0.817--0.852) against a
visit-level 0.839 --- but the distribution was strongly bimodal rather than
concentrated near that mean. Of 898 participants, 572 (63.7\%) were correct at
every eligible visit and 45 (5.0\%) at none; the median per-participant accuracy
was 1.000.

The dominant source of variation was whether a participant's recorded label
changed at all. Participants whose eligible outcomes were all TD averaged 0.944,
and those with an even mixture of TD and PIGD outcomes averaged 0.468.
\textbf{The persistence rule showed almost the same gradient on the same groups
--- 0.970 and 0.480 --- so this describes the label sequence rather than model
behaviour.} Per-participant PIGD sensitivity averaged 0.722 against a TD
specificity of 0.917 among participants with at least two outcomes of the
relevant class.

These groupings were defined after the predictions were generated and from
recorded labels only; the analysis is descriptive and is reported as
exploratory.

% ===================================================================
\section{Discussion}
% ===================================================================

\subsection{Principal findings}
Nine model specifications drawn from five dissimilar families predicted the next
recorded motor subtype within 1.22 accuracy points of one another, and
essentially all of their shared advantage over a parameter-free rule arose on
the 9.9\% of assessments at which the current label was Indeterminate. Over the
same transitions, a third of observed subtype changes were reversed at the
following visit. Read together, these results locate the limit on longitudinal
subtype prediction in this cohort in the reliability of the label rather than in
the capacity of the models.

The convergence carries more weight than any one accuracy figure. Markov chains
in the observed label, hidden Markov models with a misclassification
matrix~\cite{jackson2011}, a hierarchical subscore mixed model, Gaussian
processes on the subscores~\cite{rasmussen2006} and a time-varying spline model
embody different assumptions about what generates the data; agreement among them
is therefore evidence about the data rather than about a modelling choice.

The dissociation sharpens it. The continuous nine-domain representation has
access to far more information than the two subscores the label is computed
from, and on the two-class label task it was still statistically
indistinguishable from the hidden Markov model --- yet the same representation
discriminated later clinical outcomes better than the label did. A model does
not reach a ceiling on one target while exceeding the label on another because
it lacks information. The more parsimonious reading is that the label itself is
close to the limit of what is predictable about it, which follows from its
construction: the subtype is a deterministic threshold on a ratio of two means,
so predicting it reduces to predicting where that ratio falls relative to two
cut points.

\subsection{The Indeterminate band is the central object}
The band behaves not as a clinical state but as the neutral zone of a noisy
continuous quantity. It was the least persistent state by a wide margin, it was
the origin and destination of the changes most likely to reverse, and it is
where every fitted model found its advantage --- not by predicting it, since no
fitted specification recovered it at any useful rate, but by declining to repeat
a class that is rarely repeated. Its fragility is arithmetic as much as
clinical: half of Indeterminate visits sat at a single attainable ratio value one
hundredth from a threshold, and moving the thresholds by about 0.02 relabelled
6.7\% of all labelled visits and a fifth of participants, every one of them out
of that class. The emission constraint in our latent-state specification was
binding for this state alone, which is the same fact expressed in the model's
parameters. Whether the construct identifies a distinct phenotype has been
questioned on clinical grounds~\cite{nutt2016,urso2022}; these results add a
measurement-level reason for the same caution, confined to the Indeterminate
category.

\subsection{Implications for benchmark design}
Three implications follow for anyone training or reporting a model on a derived
clinical endpoint of this kind.

\emph{First, report a parameter-free baseline on the identical case list.} A rule
that names the current label reached 0.723 three-class accuracy here and
attained the highest macro-F1 of any specification we tested, because it was the
only one that ever predicted the minority class. An accuracy figure quoted
without that comparison is close to uninterpretable.

\emph{Second, decompose the margin.} An aggregate improvement of two to three
accuracy points looked like a uniform gain and was not: it was concentrated
entirely in one tenth of the data, and in the remaining nine tenths two of the
models were bit-for-bit identical to the trivial rule. Stratifying a reported
margin by the state at the prediction point is cheap and changes the
interpretation.

\emph{Third, measure the endpoint before modelling it.} Reversion rate,
interval dependence and threshold sensitivity are all computable without fitting
anything, and in this cohort each of them bounded what any model could achieve.
More generally, these results support assessing the longitudinal stability of a
subtype system before adopting it as an endpoint, which the field-level
appraisal found had not been done for most such systems~\cite{mestre2021}.

A fourth implication is specific to studies using this label: a conversion
endpoint defined on a single assessment overstated incidence in this cohort by
roughly a factor of two relative to one requiring a sustained change, so the
definition belongs in the protocol rather than in the analysis. Likewise, a
pooled transition matrix is an attendance-weighted average of intervals that
behave differently and is not a property of the disease.

\subsection{Comparison with prior work}
Our estimates sit coherently within the published literature where the cohorts
are comparable. The baseline distribution reproduced an independent analysis of
the same cohort to within a percentage point in each
class~\cite{simuni2016}. Prior work has measured net drift between two
assessments; we measured immediate reversion, and the two are not
interchangeable --- the same cohort yields a five-year cumulative incidence of
conversion to PIGD of 35.3\% when a single assessment defines the event and
15.3\% when the change must hold across two consecutive assessments. Lower
stability estimates reported in clinic-based
cohorts~\cite{ren2020,ren2021,kohat2021} are consistent with a more advanced
case mix and with documented disagreement between competing subtyping
algorithms. The constructive direction --- that a continuous description of the
same examination carries prognostic information the categorical label does not
--- is consistent with the appraisal that found the evidence supporting existing
subtype systems thin~\cite{mestre2021,marras2024}, and with prior predictive
work in this cohort having largely targeted latent or data-driven states rather
than the threshold rule~\cite{severson2021,fereshtehnejad2017}.

\subsection{Limitations}
This is a single-cohort study and none of it has been replicated externally; the
comparisons above are with published figures, not with analyses we performed on
another cohort. The target is a recomputation of a published rule, not an
adjudicated clinical diagnosis, so the reversion we measure is reversion of a
rule's output and includes both genuine phenotypic fluctuation and rating
variability, which this design cannot separate. Analyses were complete-case and
label loss is not random: 89.6\% of unlabelled visits lacked an OFF-medication
examination, which concentrates the loss in treated participants later in
follow-up. The ratio is undefined when the PIGD score is zero, at 16.3\% of
labelled visits; specifications that model the ratio drop those visits while
those that model the subscores retain them. Every two-class result excludes the
Indeterminate outcomes, and the resulting rise in accuracy reflects the removal
of an outcome the models never predict rather than better prediction.

All transition probabilities are conditional on attending the next visit, and
every model is given the interval to that visit, so none addresses when a
patient will next be seen. Bounded rating scales were modelled as Gaussian; the
information criteria we report assume independent observations and so favour the
latent layer more than clustered data warrant; multiplicity is largely
uncorrected across the specifications compared. The comparison of the
representation against the hidden Markov model by AUC rests on two different
validation schemes, leave-one-participant-out against five-fold
cross-validation, which favours the hidden Markov model, and the representation
additionally sees the target visit's prescribed dose. The outcome coordinates in
Section~\ref{sec:profile} were selected on the same data against which they were
compared, making each reported difference an upper bound, and the investigator
cognitive diagnosis is not fully independent of the feature set. Medication
enters all of these models as an input measured in routine care and confounded
by indication; no coefficient on it supports a causal reading. The
per-participant analysis and its grouping by label composition were specified
after the predictions were generated and are reported as descriptive. Finally,
the five-state hidden Markov specification failed to converge in the majority of
folds and is reported as a negative result about model capacity, not as a
comparable accuracy estimate.

% ===================================================================
\section{Conclusion}
% ===================================================================

Within the limits above, the measurement rather than the model was the binding
constraint on predicting the TD/PIGD motor subtype in this cohort. Five
dissimilar model families converged to within 1.22 accuracy points, their shared
advantage over a trivial rule was confined to the one tenth of assessments
recorded Indeterminate, a third of observed subtype changes reversed at the next
visit, and a threshold perturbation smaller than the measurement's own
resolution relabelled a fifth of participants.

The practical consequence is not that the TD/PIGD distinction should be
discarded, which these results do not support, but that its use as a
longitudinal endpoint should rest on a sustained-change definition, an
interval-aware treatment of follow-up, and a reported comparison against the
trivial rule --- and that where the question is progression rather than
classification, a continuous description of the same examination is the more
promising object of study. For the wider problem of learning from derived
clinical endpoints, the protocol matters more than the architecture:
harmonise the decision rule, report a parameter-free baseline on the identical
case list, decompose the margin by stratum, and measure the endpoint's own
reliability before concluding anything about a model's.

% ===================================================================
\section*{Acknowledgment}
% ===================================================================

Data used in the preparation of this article were obtained from the Parkinson's
Progression Markers Initiative (PPMI) database
(\url{www.ppmi-info.org/access-data-specimens/download-data}). For up-to-date
information on the study, visit \url{www.ppmi-info.org}. PPMI --- a public--private
partnership --- is funded by the Michael J. Fox Foundation for Parkinson's
Research and funding partners.
\textbf{[To be completed before submission: the full list of PPMI funding
partners, inserted verbatim from the current PPMI publication policy, is
required by the PPMI Data Use Agreement.]}

\section*{Data Availability}
PPMI data are available to qualified researchers from the LONI Image and Data
Archive under the PPMI Data Use Agreement and cannot be redistributed by the
authors. All analysis code, the configuration that fixes every parameter of the
pipeline, the figure- and table-generation script and the number-verification
script for this manuscript are available in the project repository.

\section*{Declarations}
\textbf{[To be completed before submission.]} Ethics approval and participant
consent (PPMI is conducted under IRB approval at each participating site and all
participants gave written informed consent; the authors should confirm the
wording required by the target journal), competing interests, funding, the
corresponding author's contact details, ORCID identifiers and author
contributions are required by most journals and are not recorded in the project
repository.

\bibliographystyle{IEEEtran}
\bibliography{refs}

\end{document}
